\section{Limitations and Future Work}

\subsection{Current Limitations}

The main limitation of the present benchmark is that the working subspace is
not the full FCI space for the entire system ladder. Instead, the Hamiltonian,
reference energy, and benchmark basis states are all defined in a
TrimCI-selected determinant space. The reduced space reproduces the exact FCI
result very well on the small calibration systems, but the larger-system
conclusions in this report should still be interpreted as reduced-space
results rather than exact full-space statements.

A second limitation is that the resource comparison is not yet fully symmetric
between the two basis families. In the current implementation, Quantum Krylov basis
states are generated from Trotterized Hamiltonian evolution, whereas the UCCSD
side is represented by selected single-generator states ranked by MP2 or CCSD
amplitudes. This is a meaningful practical comparison, but it does not exhaust
all possible hardware-aware variants of either family. In particular, a
truncated or selected Quantum Krylov construction could make the resource comparison
more balanced against selected UCCSD.

\subsection{Future Work}

One natural extension is to make the hardware-level benchmark more symmetric.
This includes repeating the resource analysis with truncated-Hamiltonian Quantum Krylov
constructions and comparing those results directly against warm-start and
selected-UCCSD variants. Such a study would clarify whether lower-cost Quantum Krylov
states can preserve enough of the expressivity advantage seen in the present
benchmark. A closely related direction is to investigate whether qubitization-
based implementations of Quantum Krylov time evolution can reduce the resource cost of
the Quantum Krylov basis relative to the current Trotterized construction.

A second direction is to expand the choice of reference state beyond
Hartree-Fock. For example, both Quantum Krylov and UCCSD could be initialized from a
stronger warm-start state such as the TrimCI reduced-space ground state rather
than the Hartree-Fock determinant. That would help separate the effect of basis
generation from the effect of the initial reference itself.

It would also be valuable to connect the present resource estimates more
directly to concrete hardware specifications. Mapping the benchmark circuits to
specific platform assumptions, such as superconducting or trapped-ion devices,
would make it possible to assess which basis family is closer to realistic
implementation on a given hardware stack.

Finally, the present comparison suggests several algorithmic extensions. One
possibility is to design a new hybrid basis family that combines the compact
expressivity of Quantum Krylov with the shallower state-preparation structure of
selected UCCSD. Another is to investigate related subspace constructions and
variations of UCC/Quantum Krylov ideas, for example Q-SENSE
\cite{patel2025qsense}. A complementary theoretical direction is to study more
carefully why Quantum Krylov bases so effectively describe the ground state in these
electronic-structure problems, and whether that behavior can be predicted from
spectral or correlation properties of the Hamiltonian.
